\section{Uniform critical recovery and the positive short trace}
\label{sec:v100-recovery}
The capacity obstruction and the positive short-trace obstruction
are different. The former has now been discharged uniformly in the
radius. We record exactly which essential-height conclusions this
supplies, and keep the latter source explicit.

Use the original residual partition from Section~\ref{sec:v98-critical}:
$d=2qr$, $\mathsf S=Y\cap\{\ell\le3d\}$ and
$\mathsf C=Y\cap\{\ell>3d\}$. For a Borel target selector $H$
supported on segments of total length at most $\delta$, write
$R_H^\varepsilon$ for its original weighted source and
$T_H^\varepsilon$ for the same-segment minimum of the two weighted
physical endpoint traces. No selector trace is introduced.

\begin{theorem}[Radius-uniform recovery without a capacity premise]
\label{thm:v100-recovery}
For the entire prescribed radius interval, the critical target source
satisfies
\begin{equation}\label{eq:v100-critical-height}
 \limsup_m\sup_{R,n,k}m^2\|\mathsf C^\varepsilon\|_\infty
          \le C d(a^2+2\Lambda_\circ d)^{1/256}.
\end{equation}
Every previously fixed Borel selection on short segments satisfies
\begin{equation}\label{eq:v100-paired-height}
 \limsup_m\sup_{R,n,k}m^2
       \|R_H^\varepsilon-T_H^\varepsilon\|_\infty
                         \le C(1+\varepsilon^{-1})\delta.
\end{equation}
In particular this bounds the complete selected source on all
segments with an actual grazing endpoint, since their paired minimum
is zero. Neither estimate requires a positive initial-gradient or
internal-incidence cutoff. The supremum is over the original exact
labels, not a neighboring-label enlargement.
\end{theorem}
\begin{proof}
The finite-count target certificate in
Theorem~\ref{thm:v98-critical-certificate} reads
\[
 \mathsf C^\varepsilon(t)\le
 C_0\int_t^{t+d}w_{1,b}(u)\,du,
 \qquad b=(a^2+2\Lambda_\circ d)^{1/2}.
\]
Its low-clock argument forces a loss of at least one half of the
initial angular density inside a full surviving forward interval.
The target-direction map is injective at each intermediate roof.
Apply \eqref{eq:v100-uncut-limit} with interval length $d$;
$b^{1/128}=(a^2+2\Lambda_\circ d)^{1/256}$ gives
\eqref{eq:v100-critical-height}.

For the other assertion the inherited same-segment trace comparison
is already a finite-count inequality on common kernels:
\[
 |R_H^\varepsilon(t)-T_H^\varepsilon(t)|
 \le C(1+\varepsilon^{-1})
              \int_{t-\delta}^{t+\delta}v^0(u)\,du.
\]
The target selector is held at $t$ and not differentiated. The Borel
segment bijection supplies multiplicity one for all joining intervals.
Equation~\eqref{eq:v100-v0} now pays the right side uniformly in the
radius, proving \eqref{eq:v100-paired-height}. A grazing endpoint has
zero original optical trace. A backward critical endpoint need not,
and is not discarded by this statement.
\end{proof}

\begin{theorem}[The original uniform pointwise reduction]
\label{thm:v100-endpoint}
Keep $q=1/8$, $\varepsilon(B)=A_0B^{-1/12}$,
$r=c\varepsilon(B)^4$, $a=\sqrt{qr}$ and $\delta=6qr$.
For all sufficiently large fixed $B$,
\begin{align}
 \mathcal E_M\le{}&\mathcal H_M(T_{\mathsf S}^{\varepsilon(B)})
       +C_MB^{-1/192}
       +C(qr)^{257/256}
       +C(1+\varepsilon(B)^{-1})qr.
                                               \label{eq:v100-endpoint}
\end{align}
Unlike the earlier family-wide reduction, this assertion has no
additional condition (SC) or finite-block margin as a hypothesis.
It uses the inherited physical and four-record analytic inputs and
the uniform pressure theorem proved above.
\end{theorem}
\begin{proof}
The running-compression reduction is an identity and upper comparison
on the original source before its roof supremum. Its residual is the
disjoint sum $\mathsf C^\varepsilon+\mathsf S^\varepsilon$.
Apply Theorem~\ref{thm:v100-recovery} to each part and the retained
positive scalar remainder. With $a^2=qr$ and $d=2qr$, the critical
cost in \eqref{eq:v100-critical-height} is $C(qr)^{257/256}$.
The paired comparison cost is $C(1+\varepsilon^{-1})qr$.
All auxiliary bands are fixed before the collision limit and enlarged
afterwards. The common uniform constants permit the radius supremum
inside that limit; no fixed-radius conclusion is used in its place.
This gives \eqref{eq:v100-endpoint}.

Under the retained outer schedule the explicit costs have sizes
$B^{-1/192}$, $O(\varepsilon^{257/64})$ and $O(\varepsilon^3)$.
They tend to zero with $B$ last. Thus decay of the positive paired
height along this schedule suffices to finish the original uniform
upper half, and hence the full pointwise theorem with the retained
lower half. That paired-height decay is not proved here.
\end{proof}

\begin{remark}[The remaining positive source]\label{rem:v100-positive}
The minimum in $T_{\mathsf S}^\varepsilon$ is taken before the
physical sheets are summed. It is not the Jordan positive part of
an already projected signed current. A small difference between the
two traces cannot bound their positive minimum. On nongrazing short
segments both traces may be positive, including at a backward
critical entrance. Theorems~\ref{thm:v100-capacity} and
\ref{thm:v100-recovery} remove the separate family-uniform capacity
premise; they do not supply a boundary deconcentration theorem for
this positive minimum. Accordingly the unrestricted upper pointwise
law, full raw local limit and unrestricted point-conditioned
consequences are not marked proved. The original target, all zero
arithmetic classes and the scalar-first normalization order remain
unchanged.
\end{remark}
